\infographicchapter{Optoacoustic Tomography (OAT)}{oat.png}
\label{ch:infographic-oat}

\section{Overview}

Optoacoustic tomography (OAT)---also widely called photoacoustic tomography---is a newer hybrid scan that combines harmless laser light with sound waves. When very brief pulses of near-infrared light shine into the body, natural substances like blood pigment (hemoglobin) absorb a tiny bit of that light energy. That tiny burst of energy makes the tissue expand slightly and send out gentle, high-pitched sound waves. Because sound waves travel through the body much more cleanly than light does, an ultrasound sensor can catch those echoes and build sharp pictures up to a few centimeters beneath the skin, without any ionizing radiation.

An advanced version called multispectral optoacoustic tomography (MSOT) shines several different wavelengths (colors) of near-infrared light one after another. By comparing how the tissue responds to each color, computer software can map oxygen-rich blood, oxygen-poor blood, fat, water, and skin pigment (melanin) without injecting any contrast dye. It can also track safe, approved medical dyes like indocyanine green (ICG) when needed to map lymph nodes or blood flow.

Clinically, OAT is still mostly a research technology, though it has reached a few important milestones. One dedicated dual optoacoustic-ultrasound breast scanner has U.S. FDA approval to work alongside standard ultrasound when checking suspicious breast lumps, and certain MSOT systems hold European CE clearance. Researchers are actively testing OAT for inflammatory bowel disease, melanoma lymph node staging, blood-vessel plaque checks, and skin tumors, while working to improve how deep it can see and standardize results across clinics.

\section{Underlying Physical Principles}
\begin{itemize}
  \item \textbf{Ultra-Short Pulses of Near-Infrared Light:} A tunable laser sends harmless, billionth-of-a-second pulses of near-infrared light into the skin, staying strictly within safety limits.
  \item \textbf{Light Turns into Sound (the Optoacoustic Effect):} Natural light-absorbing substances in the body---especially blood hemoglobin, which absorbs near-infrared light far more strongly than surrounding tissue, along with melanin, fat, and water---soak up each ultra-short pulse faster than heat can spread away. That tiny fraction-of-a-degree warming triggers a brief ripple (called thermoelastic expansion) that sends out gentle ultrasound waves.
  \item \textbf{Sharp Detail Deeper Than Light Alone:} Light scatters quickly inside soft tissue, making regular optical pictures blurry after a millimeter or two. Sound waves scatter much less, so OAT keeps sub-millimeter sharpness down to about 2 to 4~centimeters deep.
  \item \textbf{Paired Ultrasound Detection:} Standard-style ultrasound sensors pressed against the skin with cool gel listen for those returning sound waves, often alternating with regular grayscale ultrasound in the very same probe.
  \item \textbf{Multi-Color Mapping (MSOT):} By firing several wavelengths of light and sorting the echoes mathematically (a step called spectral unmixing), software separates oxygen-rich blood, oxygen-poor blood, melanin, fat, and approved dyes into color-coded maps.
\end{itemize}

\section{Imaging Mechanism}
\begin{enumerate}
  \item \textbf{Eye Protection and Gel Setup:} Everyone in the room puts on special laser safety glasses. The clinician applies clear, bubble-free ultrasound gel (or places a soft water-filled pad) on the skin and rests the handheld probe gently over the area.
  \item \textbf{Alternating Light Pulses and Ultrasound:} The laser rapidly flashes several near-infrared wavelengths into the tissue while the probe also sends out regular ultrasound pulses as the clinician glides it smoothly across the skin.
  \item \textbf{Listening for the Sound Waves:} Sensors inside the probe pick up the gentle sound waves produced by each light flash and convert them into fast digital signals.
  \item \textbf{Building Color Maps:} Software reconstructs where light was absorbed and separates out oxygen-rich hemoglobin, oxygen-poor hemoglobin, melanin, fat, or injected dye.
  \item \textbf{Live Side-by-Side Display:} Color maps of blood vessels and oxygen levels appear right on top of the live grayscale ultrasound picture, helping doctors score breast lumps or review MSOT study data immediately.
\end{enumerate}

\section{Key Technical Parameters}
\noindent
\begin{tabularx}{\textwidth}{@{}l X@{}}
\toprule
\textbf{Parameter} & \textbf{Typical Value / Description} \\
\midrule
Physical Principle & Short pulses of near-infrared light are absorbed by blood and tissue, creating gentle ultrasound waves that a sensor turns into detailed maps \\
Energy / Signal Source & Tunable pulsed near-infrared laser paired with an ultrasound probe (often combined in one handheld unit) \\
Typical Frequency / Energy & 680--980~nm near-infrared light pulses (kept within ANSI/IEC safety limits); 2--15~MHz detected sound waves \\
Spatial Resolution & About 0.1--0.5~mm detail at depths of 2--4~cm beneath the skin \\
Typical Scan Duration & 15--30~minutes for a targeted handheld breast or MSOT exam \\
Contrast Agent & None needed for natural blood, oxygen, fat, and pigment mapping; optional FDA-approved dye (indocyanine green) in select studies \\
Ionizing Radiation & None (uses non-ionizing near-infrared light and sound waves) \\
\bottomrule
\end{tabularx}

\section{Clinical Usages}

\subsection{Primary Usages}
\begin{itemize}
  \item \textbf{Checking Suspicious Breast Lumps (FDA-Cleared Adjunct):} An FDA-approved optoacoustic-ultrasound breast system works alongside standard ultrasound to map blood vessels and oxygen inside borderline (BI-RADS~3--4) breast lumps. In a multicenter clinical trial of more than 2,000~women, adding optoacoustic imaging improved diagnostic specificity by 14.9\% while keeping sensitivity high (98.6\% versus 96.0\% for ultrasound alone), helping doctors spot benign lumps and avoid unnecessary biopsies---even in dense breast tissue.
  \item \textbf{Checking Bowel Inflammation (Clinical Trials):} Handheld MSOT probes are being tested to measure blood flow in the intestinal wall without radiation or scopes, offering a gentle way to track Crohn's disease and ulcerative colitis activity.
  \item \textbf{Checking Melanoma Lymph Nodes (Research Target):} Because melanoma cells contain melanin pigment, researchers are testing MSOT---sometimes with indocyanine green dye---to check nearby sentinel lymph nodes for melanoma spread without cutting.
  \item \textbf{Evaluating Blood-Vessel Health (Clinical Research):} MSOT is being studied in poor leg circulation (peripheral arterial disease), diabetic foot care, and neck artery plaques to map tissue oxygen and fragile blood-vessel buildup.
\end{itemize}

\subsection{Secondary Usages}
\begin{itemize}
  \item \textbf{Mapping Skin Tumors (Investigational):} Pilot studies test OAT to measure the depth and borders of skin cancers like basal cell carcinoma and melanoma before surgery.
  \item \textbf{Evaluating Thyroid Nodules (Clinical Research):} Research protocols look at blood flow and oxygen inside thyroid nodules to help tell harmless nodules from cancerous ones alongside standard ultrasound.
  \item \textbf{Muscle and Joint Research (Investigational):} Clinical studies are exploring collagen, fat, and blood signals in muscle disorders (like Duchenne muscular dystrophy), skin-thickening disease (systemic sclerosis), and inflammatory arthritis.
  \item \textbf{Surgical, Biopsy, and Endocavity Guidance (Clinical Research):} Researchers are testing handheld and endocavity (transrectal or transvaginal) OAT probes to guide needle biopsies, check tumor borders during cancer surgery, protect nearby nerves during prostate operations, and evaluate deep pelvic organs such as the prostate and ovaries.
  \item \textbf{Preclinical Lab Studies:} In animal research, OAT is widely used to study tumor blood-vessel growth and how new medicines travel through the body.
\end{itemize}

\section{Advantages}
\begin{itemize}
  \item \textbf{Rich Light Contrast at Useful Depths:} Turning light absorption into sound avoids the heavy blurriness of regular light scans, keeping sharp detail down to 2 to 4~centimeters.
  \item \textbf{Dye-Free Blood and Oxygen Mapping:} MSOT maps oxygen-rich blood, oxygen-poor blood, melanin, fat, and water naturally without needing contrast injections.
  \item \textbf{No Ionizing Radiation:} Harmless near-infrared light and sound waves make OAT safe to repeat for tracking chronic inflammation or cancer treatment.
  \item \textbf{Matches Directly with Live Ultrasound:} Combining both tools in one handheld probe lays color oxygen and blood-vessel maps right over familiar grayscale ultrasound pictures.
  \item \textbf{Works with Approved Dyes When Needed:} OAT easily detects safe, approved dyes like indocyanine green for lymph-node and circulation studies.
\end{itemize}

\section{Disadvantages}
\begin{itemize}
  \item \textbf{Limited Depth:} Light fades quickly inside the body, so OAT can only see about 2 to 4~centimeters deep---not deep abdominal organs or through the adult skull.
  \item \textbf{Color Distortion at Depth:} Because overlying tissue absorbs some light colors faster than others (a problem called spectral coloring), deep oxygen measurements can become less accurate.
  \item \textbf{Bone, Air, and Darker Skin Effects:} Bone and air-filled lungs or bowel block sound waves, while darker skin tones have more natural melanin that absorbs light at the surface, reducing depth and adding extra echoes.
  \item \textbf{Laser Eye Safety and Gel Needs:} Pulsed lasers require protective goggles for everyone in the room, strict skin-energy limits, and clean acoustic gel.
  \item \textbf{Mostly Still in Research:} Aside from one FDA-approved adjunctive breast system and select CE-marked units, wider clinical use awaits standardized calibration, common rules, and large multi-hospital trials.
\end{itemize}

\section{Contraindications}

\subsection{Absolute}
\begin{itemize}
  \item \textbf{No Eye Protection or Unapproved Substitution:} Operating the laser without proper protective goggles is strictly forbidden. Outside FDA-approved adjunctive breast imaging, OAT/MSOT must never replace standard medical imaging or biopsy.
  \item \textbf{Severe Dye Allergy (If Dye Is Used):} If indocyanine green (ICG) dye is part of a study, anyone with a prior severe allergy to ICG or iodine-containing dyes cannot receive the injection.
\end{itemize}

\subsection{Relative}
\begin{itemize}
  \item \textbf{Open Skin Sores or Dark Tattoos:} Open wounds, active skin infections, bad sunburns, or dark tattoos right over the scan spot can soak up too much laser light, causing discomfort or confusing echoes.
  \item \textbf{Deep Targets, Bone, Air, or Very Dark Skin:} Targets deeper than 3 to 4~centimeters, areas hidden behind bone or air-filled bowel, or very high skin melanin can reduce depth and make color oxygen maps less reliable.
\end{itemize}

\section{Patient Safety Guidelines}
\begin{itemize}
  \item \textbf{Strict Laser Energy and Skin-Tone Safety:} Lasers are kept strictly within ANSI~Z136 and IEC~60825 safety limits for skin and eyes. Operators watch energy levels closely on darker skin tones, where extra surface melanin absorbs more light.
  \item \textbf{Mandatory Laser Goggles and Room Controls:} You, the operator, and everyone in the room must wear wavelength-matched laser safety goggles, backed up by door locks, beam shutters, and warning signs.
  \item \textbf{Clean Probes and Gel:} Staff use clean, clear ultrasound gel or water pads, disinfect handheld probes thoroughly between patients, and use sterile covers if scanning near broken skin or in surgery.
  \item \textbf{Contrast Dye Readiness:} If indocyanine green is given, staff check allergy history beforehand, confirm a clean IV line, and keep allergy-treatment medicines ready.
  \item \textbf{Clear Clinical and Research Rules:} Routine clinical use is limited to approved indications (like adjunctive breast scanning); all other OAT or MSOT human scans take place under approved research (IRB) protocols with written consent.
\end{itemize}

\section{Observing for Abnormal Patterns}
\begin{itemize}
  \item \textbf{Extra Blood Vessels and Low Oxygen in Breast Cancers:} In FDA-cleared adjunctive breast scans, cancerous lumps tend to show lots of tangled blood vessels inside and around their edges, along with lower oxygen levels (more oxygen-poor blood). Harmless fibroadenomas and cysts usually show very little internal blood flow or normal oxygen-rich flow around the rim.
  \item \textbf{High Blood Signal in Inflamed Bowel:} In Crohn's disease and ulcerative colitis MSOT studies, actively inflamed bowel segments show thicker walls packed with extra blood signal, matching inflammation seen on colonoscopy.
  \item \textbf{Melanin Spots in Lymph Nodes:} In melanoma research studies, bright melanin signals inside nearby sentinel lymph nodes---often paired with green dye tracking---are studied as signs of melanoma spread.
  \item \textbf{Depth-Related Color Shifts:} Because surface tissues absorb different colors of light unevenly (spectral coloring), deep oxygen readings can sometimes look artificially low and mimic true low oxygen.
  \item \textbf{Surface Echoes from Dark Pigment:} Strong light absorption in darkly pigmented skin can create loud surface sound waves that bounce off bone, air, or blood vessels, leaving faint ``ghost'' echoes deeper in the image.
\end{itemize}

\section{Dealing with Uncertainties}
\begin{itemize}
  \item \textbf{Light Fading and Color Distortion:} While laboratory test models can show signals 5 to 6~centimeters deep, real-world clinical scans typically reach only 2 to 3~centimeters (or up to 4~centimeters in breast tissue). As light weakens rapidly with depth, deeper regions can look artificially dim---or show skewed oxygen numbers---unless software corrects for uneven light loss.
  \item \textbf{Skin-Tone and Echo Clutter:} Higher melanin in darker skin absorbs more surface light, limiting how deep the scan sees and creating bouncing echo clutter.
  \item \textbf{Probe Angle and Movement Shifts:} Handheld probes view tissue from only one side, which can slightly stretch shapes, while breathing or pressing harder with the probe can shift the color map away from the grayscale ultrasound.
  \item \textbf{Active Benign Lumps Can Overlap:} Some harmless lumps (like inflamed fibroadenomas, papillomas, or swollen reactive lymph nodes) also have extra blood flow, creating borderline scores.
  \item \textbf{Confirming with Proven Care:} Operators cut down artifacts by using bubble-free gel, steady light pressure, calibration checks, and multiple angles. Any uncertain breast lump follows standard BI-RADS ultrasound, mammogram, MRI, or needle biopsy care, while research MSOT results are confirmed by colonoscopy or surgical pathology.
\end{itemize}

\section{Augmenting Methods}
\begin{itemize}
  \item \textbf{Paired Grayscale and Doppler Ultrasound:} Standard ultrasound is built right into handheld OAT probes, giving clear anatomical borders and larger blood-vessel flow alongside OAT's tiny-vessel oxygen maps.
  \item \textbf{Mammography and Breast MRI:} Mammograms catch tiny calcium specks that OAT cannot see, while contrast-enhanced breast MRI sees deep into the whole breast beyond OAT's 2 to 4~centimeter depth limit.
  \item \textbf{Colonoscopy and CT/MR Enterography:} In inflammatory bowel disease studies, colonoscopy with tissue biopsy and full-abdomen MRI or CT show deep bowel loops that handheld OAT cannot reach.
  \item \textbf{Needle Biopsy and Pathology:} Ultrasound-guided core-needle biopsy for suspicious breast lumps and lab pathology of lymph nodes remain the final gold standards for cancer diagnosis.
  \item \textbf{Smarter Light-Correction Software:} New computer models that adjust for depth, skin tone, and tissue phantoms are helping make oxygen measurements more accurate.
\end{itemize}

\section{Comparison with Alternative Modalities}
\begin{itemize}
  \item \textbf{Versus Standard Ultrasound:} OAT adds dye-free blood-oxygen and tiny-vessel detail right onto standard ultrasound, while standard ultrasound sees deeper, costs less, needs no laser goggles, and is available everywhere.
  \item \textbf{Versus Breast MRI:} Breast MRI sees the entire depth of both breasts for full screening and staging, while adjunctive optoacoustic-ultrasound is faster, far less costly, needs no gadolinium injection, and works right at the bedside for checking borderline lumps.
  \item \textbf{Versus Pure Optical Scans and Photoacoustic Microscopy:} Turning light into sound lets wide-field OAT see centimeters deep instead of just a millimeter or two. Focused point-by-point variants (photoacoustic microscopy) trade that extra depth for even finer, micron-scale maps of tiny surface vessels, while pure optical microscopes see individual cells right at the surface.
  \item \textbf{Versus Mammography:} OAT uses no X-rays, avoids breast squeezing, and works well in dense breasts, while mammography remains essential for whole-breast screening and spotting microcalcifications.
\end{itemize}

\section{Key References}
\begin{itemize}
  \item Neuschler EI, Butler R, Young CA, et al. A pivotal study of optoacoustic imaging to diagnose benign and malignant breast masses: a new evaluation tool for radiologists. \textit{Radiology}. 2018;287(2):398--412.
  \item Ntziachristos V, Razansky D. Molecular imaging by means of multispectral optoacoustic tomography (MSOT). \textit{Chemical Reviews}. 2010;110(5):2783--2794.
  \item Knieling F, Neufert C, Hartmann A, et al. Multispectral optoacoustic tomography for assessment of Crohn's disease activity. \textit{New England Journal of Medicine}. 2017;376(13):1292--1294.
  \item American National Standards Institute (ANSI). \textit{ANSI Z136.1: American National Standard for Safe Use of Lasers}. Orlando, FL: Laser Institute of America, 2022.
  \item Steinberg I, Huland DM, Vermesh O, Frostig HE, Tummers WS, Gambhir SS. \textit{Photoacoustic clinical imaging}. Stanford University School of Medicine (NIH PMC). Available at: \url{https://pmc.ncbi.nlm.nih.gov/articles/PMC6595011/}.
\end{itemize}

\section{Further Reading}
\begin{itemize}
  \item Wikipedia contributors. \textit{Photoacoustic imaging}. Wikipedia, The Free Encyclopedia. Available at: \url{https://en.wikipedia.org/wiki/Photoacoustic_imaging}.
\end{itemize}
